\makeabstract
{
Identifying Multiple-star Systems in the Ultracool Dwarf Sample
}
{
Piper Lentz (1,2), Chris Theissen (1)
}
{
\textsuperscript{1}UC San Diego, La Jolla, CA, USA\\
\textsuperscript{2}Amherst College, Amherst, MA, USA
}
{
The lowest mass stars and brown dwarfs, also known as ultracool dwarfs (UCDs; objects less than \textasciitilde{}15\% the mass of the Sun), are a crucial population in the search for habitable worlds. However, exoplanet population statistics have yet to be well studied among UCDs because of inefficient exoplanet detection methods around these faint objects, resulting in the need for target prioritization for the most probable exoplanet hosts. A potential indicator of exoplanet host status is a spectroscopic peculiarity shared by TRAPPIST-1 and Teegarden’s star, both of which are exoplanet hosts with at least two planets in their habitable zones.
}
{
A sample of over 800 UCDs has been constructed, and their low-resolution, near-infrared spectra have been collected over the past 5 years. This sample is key to testing the potential correlation between the spectral peculiarity in the UCDs and the existence of planetary systems. However, first we must separate the multi-star systems from the single-star systems to maintain statistical robustness. We use the Gaia Archive and Washington Double Star Catalogue to identify UCDs in our sample that are part of multiple-star systems.
}
{
We found that 93 of our target UCDs are a part of multiple-star systems, 17 being triple or higher degree systems, and 76 being binary systems.
}
{
These multi-star systems are interesting targets to follow up with and test if the stars in the system also show this spectral peculiarity, therefore potentially being multi-star systems hosting exoplanets. Additionally, since \textasciitilde{}50\% of all Sun-like stars are in binary systems (a multi-star system consisting of two stars), and planets have been discovered in binary systems (circumbinary planets), these systems are essential to understand, especially for upcoming missions like the Habitable Worlds Observatory, which will look at planets in binary systems.
}

\newpage
